%%%%%%%%%%%%%%%%%%%%%%%%%%%%%%%%%
%
%       EXERCÍCIO
%
%%%%%%%%%%%%%%%%%%%%%%%%%%%%%%%%%

\definevalorquestao

\begin{exercicioBanco}[\valorquestao]
O tempo diário de uso do celular por estudantes de uma turma foi registrado, em horas, durante os últimos 40 dias. Os valores observados são apresentados abaixo.

\begin{center}
    \begin{tabular}{|C{30pt}|C{30pt}|C{30pt}|C{30pt}|C{30pt}|C{30pt}|C{30pt}|C{30pt}|C{30pt}|C{30pt}|}
        \hline
        2.1 & 2.2 & 2.3 & 2.5 & 2.7 & 2.9 & 3.0 & 3.0 & 3.1 & 3.1 \\
        \hline
        3.1 & 3.2 & 3.2 & 3.2 & 3.3 & 3.3 & 3.3 & 3.3 & 3.4 & 3.4 \\
        \hline
        3.4 & 3.4 & 3.4 & 3.5 & 3.5 & 3.5 & 3.6 & 3.6 & 3.6 & 3.7 \\
        \hline
        3.7 & 3.7 & 3.8 & 3.9 & 4.0 & 4.1 & 4.1 & 4.2 & 4.3 & 4.4 \\
        \hline
    \end{tabular}
\end{center}

\begin{enumerate}[(a)]
    \item Construa o histograma de tal variável estatística, considerando o comprimento de cada faixa igual a \(0.4\), com a primeira faixa começando em \(2\).
    %
    \item Construa o boxplot de tal variável estatística.
\end{enumerate}
\end{exercicioBanco}


